\title{Direct Preference Optimization: Your Language Model is Secretly a Reward Model}

\begin{abstract}
Although large-scale unsupervised language models (LMs) are capable of acquiring extensive world knowledge and demonstrating certain reasoning abilities, exerting fine-grained control over their outputs remains a challenge due to their fully unsupervised training regimes. To address this limitation, prior work typically involves collecting human feedback in the form of pairwise preferences over generated outputs and subsequently fine-tuning the LM to better reflect these preferences, frequently through reinforcement learning from human feedback (RLHF). RLHF, however, is a multi-stage and often fragile process, beginning with the construction of a reward model that captures human judgment and followed by tuning the unsupervised LM to maximize this reward through reinforcement learning, while constraining it to remain close to the initial model. In this work, we propose a novel formulation of the reward model within RLHF that permits direct derivation of the optimal policy in closed form, which enables resolution of the standard RLHF objective by optimizing a simple classification loss. This method, termed \textit{Direct Preference Optimization} (DPO), offers improved stability, strong empirical performance, and computational efficiency, eliminating the necessity for language model sampling during fine-tuning and reducing the reliance on extensive hyperparameter searches. Empirical evaluations demonstrate that DPO fine-tunes LMs to reflect human preferences as effectively as, or better than, established approaches. In particular, DPO surpasses PPO-based RLHF in tasks requiring sentiment control, and delivers comparable or superior results for summarization and single-turn dialogue, all with markedly reduced implementation complexity and computational cost.
\end{abstract}

\section{Introduction}
Large unsupervised language models (LMs) developed through training on massive datasets have demonstrated a range of unexpected capabilities~\citep{chowdhery2022palm, brown2020language, touvron2023llama,bubeck2023sparks}. Despite their impressive abilities, these models are trained on corpora created by humans who hold diverse intentions, standards, and competencies. Consequently, not every human behavior or goal reflected in the training data is desirable for emulation. For instance, it is useful for an AI coding assistant to \textit{recognize} frequent programming errors to facilitate their correction; however, during code generation, we prefer that the model favor the infrequent but superior coding exemplars found within its training corpus. Similarly, while awareness of widely held misconceptions—such as those believed by 50\% of individuals—may be important for the model, it is not appropriate for the model to perpetuate such misconceptions in half of its outputs related to the topic. Therefore, the ability to \emph{select} specific, preferred behaviors and outputs from among the model's extensive \textit{knowledge and capabilities} is essential for constructing AI systems that are reliable, effective, and amenable to control~\citep{ouyang2022training}. While current approaches frequently employ reinforcement learning (RL) to align LM outputs with human preferences, this work demonstrates that the standard RL objective underlying these methods can, in fact, be solved exactly using a straightforward binary cross-entropy loss, significantly reducing the complexity of the preference learning process.

Broadly speaking, most techniques embed desired attributes in language models by utilizing well-designed datasets of human preferences, highlighting responses and behaviors that people deem safe and beneficial. This process of preference learning takes place subsequent to extensive unsupervised pre-training on large-scale text corpora. The simplest strategy for preference alignment involves supervised fine-tuning with exemplars of excellent human-authored responses. However, the prevailing and more successful paradigm employs reinforcement learning from human (or artificial) feedback (RLHF/RLAIF; \citep{christiano2017deep,bai2022constitutional}). In RLHF, a reward model is trained on annotated human preference data, and RL is used to adjust the LM policy so as to increase the frequency of responses with higher reward while minimizing deviation from the original pretrained model. Although RLHF enables the development of models with remarkable conversational and programming performance, its pipeline is substantially more elaborate than standard supervised learning: it requires training several LMs, as well as repeatedly sampling from the policy during training, resulting in much higher computational expenditure.

This work presents a method for optimizing language models to better reflect human preferences without relying on explicit reward modeling or reinforcement learning frameworks. We introduce \textit{{Direct Preference Optimization} (DPO)}, an approach that implicitly targets the same optimization objective as traditional RLHF methods—namely, reward maximization under a KL-divergence regularization—yet offers greater simplicity in both implementation and training procedures. Conceptually, DPO adjusts the model by boosting the log-probability ratio of preferred over non-preferred outputs, employing an adaptive, example-specific weighting scheme. This mechanism mitigates the risk of model collapse seen in straightforward probability ratio-based formulations. As with previous techniques, DPO builds upon a formal preference model, such as the Bradley-Terry framework (\cite{bradley1952rankanalysis}), to quantify the correspondence between a reward function and observed preference data. Unlike standard approaches that leverage this model first to craft a preference-based loss for reward model training, subsequently training a policy to maximize this reward, DPO exploits a reparameterization strategy: the preference loss is defined directly as a function of the policy itself. Given human preference-annotated datasets for model outputs, DPO permits policy optimization using a binary cross entropy loss, effectively producing a policy optimal with respect to an implicit reward constructed from the preference information.

The principal contribution of this study is Direct Preference Optimization (DPO), a streamlined, reinforcement learning-free algorithm for learning language models from preference data. Empirical results demonstrate that DPO matches or surpasses current methods—including PPO-driven RLHF—for tasks like sentiment control, summarization, and conversational response, with language models scaling up to 6B parameters.

\section{Related Work}

As self-supervised language models grow in scale, their ability to perform certain tasks with no prior examples (zero-shot; \citep{radford2019language}) or with minimal examples (few-shot; \citep{gpt3,megatron,chowdhery2022palm}) has improved. Nevertheless, direct fine-tuning on datasets containing explicit instructions paired with human-authored responses markedly boosts their efficacy on downstream applications and better aligns their outputs with user goals \citep{mishra-etal-2022-cross,sanh2022multitask,chung2022scaling,thoppilan2022lamda}. Through this process, referred to as `instruction-tuning', language models develop an enhanced capacity to generalize beyond the range of provided instructions and to deliver more user-friendly responses \citep{chung2022scaling}. Yet, even with the effectiveness of instruction-tuning, it is often more feasible to gather comparative judgments from humans regarding response quality than to assemble large sets of expert demonstrations. As a result, subsequent approaches have leveraged datasets built from human preferences to further fine-tune language models, resulting in notable gains for translation \citep{kreutzer-etal-2018-reliability}, summarization \citep{stiennon2022learning,ziegler2020finetuning}, narrative generation \citep{ziegler2020finetuning}, and adherence to instructions \citep{ouyang2022training,ramamurthy2023is}. These methods typically entail an initial stage in which a neural reward function is trained for alignment with human preference data—using frameworks like the Bradley-Terry model \citep{bradley1952rankanalysis}—followed by language model fine-tuning to optimize this reward via reinforcement learning strategies such as REINFORCE \citep{williams1992reinforce}, proximal policy optimization (PPO; \cite{schulman2017proximal}), or related algorithms \citep{ramamurthy2023is}. Closely aligned research utilizes instruction-following LLMs, themselves refined via human feedback, to autonomously produce synthetic preference data for traits like safety or harmlessness \citep{bai2022constitutional}, relying predominantly on high-level, rubric-style textual supervision from humans to guide the model’s annotations. Altogether, these lines of inquiry merge advances in reinforcement learning–based language model training for diverse objectives~\citep{Ranzato2015SequenceLT,paulus2018a,wu2018learning} with general frameworks for inferring optimal behaviors from human preference signals \citep{christiano2017deep,kupcsik2018learning}. However, despite the appeal of working with relative preference data, the reinforcement learning techniques commonly employed to fine-tune large language models remain formidable in terms of implementation and resource demands. The present work introduces a methodologically grounded alternative that enables direct optimization of relative preferences, circumventing the need for reinforcement learning.

Beyond applications in language, the challenge of policy learning from preference information has received attention in both bandit and reinforcement learning frameworks, with a range of methods introduced in the literature. When contextual bandit models are trained using preferences or rankings over actions as opposed to explicit reward signals, the problem is referred to as contextual dueling bandits (CDB; \cite{yue2012karmed,dudik2015contextual}). Since absolute rewards are unavailable in these scenarios, theoretical work on CDBs adopts the concept of a \textit{von Neumann winner} to characterize optimality—a policy whose expected probability of prevailing over any competing policy is no less than 50\% \citep{dudik2015contextual}. Nonetheless, a distinctive aspect of the CDB setting is that preference feedback is acquired interactively during learning, whereas in the human preference learning context, the typical assumption is access to a static dataset of offline preferences annotated over action pairs \citep{yan2022human}. In a similar vein, \textit{preference-based RL} (PbRL) utilizes binary preferences generated by an underlying, unknown `scoring' function instead of observed rewards \citep{BusaFekete2014,ruiz2023dueling}. Several PbRL algorithms are available, some of which facilitate reuse of off-policy preference datasets; these generally operate by first learning an explicit model of the hidden scoring (or reward) function before optimizing the resulting model \citep{jain2013learning,BusaFekete2014,christiano2017deep,sadigh2017active,kupcsik2018learning}. Contrasting with these multi-step methodologies, our approach seeks to directly learn a policy in a single stage, optimizing for the satisfaction of provided preference data.

\section{Preliminaries}\label{section:prelims}

We begin by summarizing the RLHF workflow as outlined by \citeauthor{ziegler2020finetuning}, and as subsequently refined by works such as \citep{stiennon2022learning, bai2022training, ouyang2022training}. This procedure typically comprises three primary components: (1) supervised fine-tuning (SFT); (2) reward model training from preference data; and (3) optimization through reinforcement learning.

\textbf{SFT}: The RLHF process customarily starts with the supervised fine-tuning of a large, pre-trained language model (LM) using curated datasets tailored to the target task (e.g., dialogue, summarization), yielding a policy denoted $\pi^\text{SFT}$.

\textbf{Reward Modelling Phase}: During the subsequent reward learning phase, prompts $x$ are provided to the SFT model, generating pairs of candidate completions $(y_1, y_2)\sim \pi^\text{SFT}(y \mid x)$. Human annotators then compare these outputs and indicate which answer they prefer, with $y_w\succ y_l \mid x$ representing that $y_w$ is favored over $y_l$ given $x$, where $y_w$ and $y_l$ correspond to the chosen and rejected responses respectively. While the precise reward mechanism, $r^*(y, x)$, underlying these choices remains unknown, various models have been employed to capture preference information, among which the Bradley-Terry (BT) model \cite{bradley1952rankanalysis} is commonly adopted. More sophisticated frameworks such as the Plackett-Luce ranking model \citep{plackett1975analysis, luce2012individual} can be leveraged for settings with multiple rankings. The BT model specifies that the true (human) preference probability $p^*$ is formalized as follows:

\begin{equation}\label{eq:bradley-terry}
    p^*(y_1\succ y_2 \mid x)=\frac{\exp\left(r^*(x, y_1)\right)}{\exp\left(r^*(x, y_1)\right) + \exp\left(r^*(x, y_2)\right)}.
\end{equation}

Given a fixed collection of preference data $\mathcal{D}=\bigl\{x^{(i)}, y_w^{(i)}, y_l^{(i)}\bigr\}_{i=1}^N$ drawn according to $p^*$, we may introduce a reward model parameterization $r_{\phi}(x, y)$ and fit its parameters by maximizing the likelihood of observed data. This setting can be interpreted as a binary classification problem, leading to the following negative log-likelihood objective:

\begin{equation}\label{eq:reward_model}
    \mathcal{L}_R(r_{\phi}, \mathcal{D}) = -\mathbb{E}_{(x, y_w, y_l)\sim \mathcal{D}}\bigl[\log \sigma(r_{\phi}(x, y_w)- r_{\phi}(x, y_l))\bigr]
\end{equation}

where $\sigma$ denotes the logistic sigmoid. Within language model frameworks, $r_{\phi}(x, y)$ is typically seeded from the SFT model $\pi^\text{SFT}(y \mid x)$, and an additional linear transformation is appended atop the last transformer block to generate a scalar-valued reward prediction \cite{ziegler2020finetuning}. Previous works have noted the benefit of normalizing the reward values such that $\mathbb{E}_{x,y\sim \mathcal{D}}\left[r_\phi(x, y)\right] = 0$ across all $x$, in order to control the variance of the learned reward signal.

\textbf{RL Fine-Tuning Phase}: During reinforcement learning fine-tuning, the policy receives feedback from the learned reward function. As described in prior work~\citep{jaques2017sequence, jaques2020human}, the objective to optimize becomes

\begin{equation}\label{eq:RL}
\max_{\pi_{\theta}}  \mathbb{E}_{x\sim \mathcal{D}, y\sim \pi_{\theta}(y \mid x)}\bigl[r_{\phi}(x, y)\bigr] - \beta\mathbb{D}_{\textrm{KL}}\bigl[\pi_{\theta}(y\mid x)\mid \mid \pi_\text{ref}(y\mid x)\bigr],
\end{equation}

with $\beta$ serving as a regularization parameter that limits the divergence from the reference policy $\pi_\text{ref}$, which is ordinarily set to the initial SFT model $\pi^\text{SFT}$. The initial parameters of $\pi_\theta$ are typically copied from $\pi^\text{SFT}$ as well. Imposing this KL penalty is critical to ensure that the policy remains close to the region where the reward estimator is reliable, in addition to promoting diverse generations and avoiding collapse into a narrow set of high-reward outputs. Because language model outputs are discrete, the objective above cannot be differentiated directly and is usually addressed using reinforcement learning techniques. Standard practice \citep{ziegler2020finetuning, stiennon2022learning, bai2022training, ouyang2022training} constructs the reward as ${r(x, y) = r_{\phi}(x, y) -\beta (\log \pi_{\theta}(y\mid x) - \log \pi_\text{ref}(y\mid x))}$ and then employs PPO \cite{schulman2017proximal} to optimize it.

\section{Direct Preference Optimization}\label{sec:DPO}

Given the considerable difficulty of applying reinforcement learning approaches at scale for fine-tuning large language models, we aim to develop a more straightforward method for policy optimization based directly on preferences. In contrast to conventional RLHF pipelines, which first fit a reward model and then carry out RL optimization, our strategy makes use of a particular form of reward parameterization that admits an analytical solution for the associated optimal policy, sidestepping the need for iterative RL training. In detail, the core idea is to leverage a closed-form correspondence between reward functions and their optimal policies, enabling the recasting of the reward modeling loss into an equivalent objective for the policy. This change-of-variable technique eliminates the requirement for an explicit separate reward model, while still being grounded in established models for capturing human preferences, such as the Bradley-Terry formulation.

\textbf{Derivation of the DPO Objective.} We begin from the standard reinforcement learning objective defined in Eq.~\ref{eq:RL}, utilizing a general reward function $r$, as has been done in previous research. Drawing on established literature~\citep{peters2007reinforcement, peng2019advantage, korbak2022reinforcement, go2023aligning}, it can be readily demonstrated that the policy that optimizes the KL-regularized reward maximization problem in Eq.~\ref{eq:RL} is characterized by:

\begin{equation}\label{eq:op_policy}
    \pi_r(y\mid x) = \frac{1}{Z(x)}\pi_\text{ref}(y\mid x)\exp\left(\frac{1}{\beta}r(x, y)\right),
\end{equation}

where the normalization constant $Z(x) =\sum_{y}\pi_\text{ref}(y\mid x)\exp\left(\frac{1}{\beta}r(x, y)\right)$ ensures the distribution sums to one. The full mathematical steps are provided in Appendix \ref{app:derivation1}. Even when employing the maximum likelihood estimate $r_\phi$ of the true reward $r^*$, determining $Z(x)$ exactly is computationally demanding~\citep{korbak2022reinforcement, go2023aligning}, limiting the practicality of this formulation. To address this, Eq.~\ref{eq:op_policy} can be re-expressed to write the reward as a function of the resulting optimal policy $\pi_r$, the reference policy $\pi_\text{ref}$, and the partition function $Z(\cdot)$. Taking the logarithm of Eq.~\ref{eq:op_policy} and performing algebraic manipulations yields:

\begin{equation}\label{eq:main_eq}
    r(x,y) =\beta \log \frac{\pi_r(y\mid x)}{\pi_\text{ref}(y\mid x)} + \beta \log Z(x).
\end{equation}

This transformation can be directly applied to the true reward $r^*$ and its associated optimal policy $\pi^*$. Importantly, the Bradley-Terry model is sensitive only to the difference in reward values between two outputs, that is, ${p^*(y_1 \succ y_2 \mid x) = \sigma(r^*(x, y_1) - r^*(x, y_2))}$. By substituting the reformulation from Eq.~\ref{eq:main_eq} for $r^*(x,y)$ in the preference model given by Eq.~\ref{eq:bradley-terry}, the normalization term cancels, and the human preference probability can thus be formulated exclusively in terms of the optimal policy $\pi^*$ and the reference policy $\pi_\text{ref}$. Consequently, under the Bradley-Terry model, the optimal RLHF policy $\pi^*$ satisfies the following preference likelihood:

\begin{equation}\label{eq:objective}
    p^*(y_1\succ y_2 \mid x)=\frac{1}{1 + \exp\left(\beta \log \frac{\pi^*(y_2\mid x)}{\pi_\text{ref}(y_2\mid x)} - \beta \log \frac{\pi^*(y_1\mid x)}{\pi_\text{ref}(y_1\mid x)}\right)}
\end{equation}

A comprehensive derivation can be found in Appendix~\ref{app:derivation2}. Although this expression employs the Bradley-Terry framework, analogous derivations are possible using the more general Plackett-Luce models~\citep{plackett1975analysis, luce2012individual}, as detailed in Appendix~\ref{app:plackett_luce_models}.

Given this formulation where human preference probabilities are parametrized by the policy rather than by the reward function, we may now pose a maximum likelihood training objective for a policy $\pi_\theta$. Drawing a parallel with the reward modeling loss (see Eq.~\ref{eq:reward_model}), we obtain the following objective for policy optimization:

\begin{equation}\label{eq:optimum_model}
    \mathcal{L}_\text{DPO}(\pi_{\theta}; \pi_\text{ref}) = -\mathbb{E}_{(x, y_w, y_l)\sim \mathcal{D}}\left[\log \sigma \left(\beta \log \frac{\pi_{\theta}(y_w\mid x)}{\pi_\text{ref}(y_w\mid x)} - \beta \

In this framework, we model an implicit reward through an alternative parameterization, with the optimal policy corresponding directly to $\pi_\theta$. Additionally, our approach is mathematically analogous to fitting a reparametrized Bradley-Terry model, which ensures certain desirable theoretical characteristics, such as consistency given appropriate assumptions about the preference data distribution \cite{bong2022generalized}. Further analysis of DPO's theoretical properties, particularly in the context of related methods, is provided in Section~\ref{sec:theory}.

\textbf{How does the DPO update operate?} To gain mechanistic insight into DPO, it is instructive to study the gradient of the objective function $\mathcal{L}_\text{DPO}$. The derivative with respect to the parameter vector $\theta$ is expressed as:

\begin{multline*}\label{eq:gradient}
    \nabla_\theta \mathcal{L}_\text{DPO}(\pi_\theta;\pi_\text{ref}) = \\ -\beta\mathbb{E}_{(x, y_w, y_l) \sim \mathcal{D}} \bigg[\underbrace{\sigma(\hat{r}_\theta(x, y_l) - \hat{r}_\theta (x, y_w))}_\text{larger impact when reward ranking is incorrect}\bigg[\underbrace{\nabla_\theta\log \pi(y_w \mid x)}_\text{promote $y_w$ occurrence} - \underbrace{\nabla_\theta\log\pi(y_l \mid x)}_\text{demote $y_l$ occurrence}\bigg]\bigg],
\end{multline*}

Here, the term $\hat{r}_\theta(x, y) = \beta \log \frac{\pi_\theta(y \mid x)}{\pi_\text{ref}(y \mid x)}$ defines the implicit reward computed via the language model $\pi_\theta$ relative to a reference $\pi_\text{ref}$ (see Section~\ref{sec:theory} for details). Conceptually, this gradient step shifts the model to favor more likely preferred responses $y_w$, while reducing the likelihood of less favored options $y_l$. Crucially, each training example is weighted by the degree to which the implicit reward function $\hat{r}_\theta$ overvalues the dispreferred $y_l$, modulated by $\beta$—that is, by the extent to which the model misorders the pair, in accordance with the KL penalty strength. Empirically, we observe that such weighting is vital; omitting this factor, as in a simplified version of the update, can cause the model’s generations to deteriorate (see Appendix Table~\ref{tab:unlikelihood_generations}).

\textbf{DPO overview.} The standard DPO process consists of the following steps: First, generate completion samples $y_1, y_2 \sim \pi_\text{ref}(\cdot \mid x)$ for each prompt $x$, and annotate these with human preferences to create the offline preference dataset $\mathcal{D} = \{x^{(i)}, y_w^{(i)}, y_l^{(i)}\}_{i=1}^N$. Next, the language model $\pi_\theta$ is trained to minimize $\mathcal{L}_\text{DPO}$, conditioned on the specified $\pi_\text{ref}$, dataset $\mathcal{D}$, and the chosen $\beta$. In practice, rather than generating new samples and collecting fresh human preference labels, researchers often leverage existing, publicly available preference datasets. Because these datasets are generally collected from completions sampled by $\pi^\text{SFT}$, the reference policy $\pi_\text{ref}$ is set to $\pi^\text{SFT}$ if it exists. Alternatively, in the absence of $\pi^\text{SFT}$, the initialization for $\pi_\text{ref}$ involves maximizing the likelihood over the preferred completions ${(x, y_w)}$, specifically: ${\pi_\text{ref} = \argmax_{\pi}\mathbb{E}_{x, y_w \sim \mathcal{D}}\left[\log \pi(y_w \mid x)\right]}$. This strategy reduces the mismatch between the unavailable true reference distribution and the $\pi_\text{ref}$ employed in DPO. Further specifics on implementation and chosen hyperparameters are discussed in Appendix~\ref{app:implementation}.

\section{Theoretical Analysis of DPO} This section provides deeper insight into the DPO approach, establishing its theoretical foundation, and contrasts DPO's strengths with the challenges associated with actor-critic algorithms commonly applied to RLHF, such as PPO~\cite{schulman2017proximal}.

\label{sec:theory}

\subsection{Your Language Model Is Secretly a Reward Model} DPO unifies reward modeling and policy learning under a single maximum likelihood framework, effectively removing the need for separate reward fitting and RL-based policy optimization steps. Importantly, the objective function in Eq. \ref{eq:main_eq} can be interpreted as a Bradley-Terry model where the rewards are given by $r^*(x, y) = \beta \log\frac{\pi^*_\theta(y \mid x)}{\pi_\text{ref}(y \mid x)}$. Here, we optimize the parameterized policy $\pi_{\theta}$ in a manner analogous to optimizing a reward model as described in Eq. \ref{eq:reward_model}, following a variable transformation. This section develops the theoretical underpinnings for this reparameterization, demonstrating that it does not limit the expressiveness of the learned reward functions and that it permits the precise recovery of the optimal policy. We start by introducing an equivalence relation among reward functions.

\begin{definition}
Two reward functions $r(x, y)$ and $r'(x, y)$ are considered equivalent if and only if there exists some function $f$ such that $r(x, y) - r'(x, y) = f(x)$.
\end{definition}

One can readily verify that this defines an equivalence relation, which in turn partitions the set of all reward functions into distinct equivalence classes. We present the following two lemmas:

\begin{lemma}\label{lemma:same_prefrence}
Within the Plackett-Luce framework, including the Bradley-Terry model as a special case, any pair of reward functions belonging to the same equivalence class induces identical preference distributions.
\end{lemma}

\begin{lemma}\label{lemma:same_policy}
    Any two reward functions in the same equivalence class yield the same optimal policy for the corresponding constrained RL formulation.
\end{lemma}

The proofs follow directly and are included in Appendix \ref{app:lemma1}. The first lemma corresponds to the well-documented under-identifiability inherent in models of the Plackett-Luce family \cite{plackett1975analysis}. As a consequence of this ambiguity, additional constraints are commonly introduced to ensure identifiability of maximum likelihood estimates in Eq. \ref{eq:reward_model} \cite{bong2022generalized}. The second lemma asserts that all reward functions from a given equivalence class lead to the same optimal policy. Thus, for our principal objective, it suffices to recover any representative reward function from the optimal class. We formalize the following Theorem and provide its proof in Appendix~\ref{app:thm1}:

\begin{theorem}\label{thm:main}
    Given mild conditions, every reward equivalence class consistent with Plackett-Luce (including Bradley-Terry) models can be expressed via the reparameterization ${r(x, y) = \beta \log \frac{\pi(y\mid x)}{\pi_\text{ref}(y\mid x)}}$ for some model $\pi(y\mid x)$ and fixed reference model $\pi_\text{ref}(y \mid x)$.
\end{theorem}

\begin{sproof}
    Let $r(x, y)$ be any reward function, associated with an optimal model $\pi_r(y \mid x)$ as specified in Eq. \ref{eq:op_policy}. We aim to show that there is a member of the equivalence class of $r$ that admits the stated reparameterized form. Define the projection $f$ by  
\begin{equation}
    f(r; \pi_\text{ref}, \beta)(x, y) = r(x, y) - \beta\log\sum_{y}\pi_\text{ref}(y\mid x)\exp\left(\frac{1}{\beta}r(x, y)\right)
\end{equation}
Here, the operator $f$ normalizes the reward function by subtracting the log-partition function of $\pi_r$, which depends solely on $x$. Therefore, $f(r; \pi_\text{ref}, \beta)(x, y)$ remains in the equivalence class of $r(x, y)$. Substituting the expression for $r$ using Eq.~\ref{eq:main_eq} (valid for any reward function), we obtain $f(r; \pi_\text{ref}, \beta)(x, y) = \beta \log \frac{\pi_r(y\mid x)}{\pi_\text{ref}(y\mid x)}$. Hence, this projection yields a member of the same equivalence class in the required form, and the proposed reparameterization encompasses the full expressivity required for our reward modeling.
\end{sproof}

Theorem~\ref{thm:main} can also be interpreted as precisely identifying the specific reward function chosen by the DPO reparameterization within each equivalence class; namely, the reward function that fulfills the condition:

\begin{equation}\label{eq:lag_p}
     \sum_{y}\underbrace{\pi_\text{ref}(y\mid x)\exp\left(\frac{1}{\beta}r(x, y)\right)}_{=\pi(y\mid x)\text{, using Thm.~\ref{thm:main} reparam.}} = 1,
\end{equation}

which ensures that $\pi(y\mid x)$ defines a valid probability distribution (non-negative and normalized to 1). By referencing Eq.~\ref{eq:op_policy}, it is evident that Eq.~\ref{eq:lag_p} serves as the partition function for the optimal policy determined by the reward function $r(x, y)$. The DPO algorithm’s central contribution lies in constraining the typically under-specified Plackett-Luce (and, specifically, the Bradley-Terry) family of preference models. Through these constraints, the full set of reward models remains representable, but the optimal policy described in Eq.~\ref{eq:op_policy} becomes analytically manageable for every prompt $x$.

\subsection{Instability of Actor-Critic Algorithms} Our proposed framework further enables an analysis of why conventional actor-critic algorithms for RLHF, such as PPO, may exhibit instability. Adopting the standard RLHF setup, we focus on the fine-tuning process as detailed in Section \ref{section:prelims}. Drawing parallels with the control as inference approach \cite{levine2018reinforcement}, applied to the constrained RL formulation presented in \ref{eq:RL}, we consider a parameterized policy model $\pi_{\theta}(y\mid x)$ and aim to minimize the divergence $\mathbb{D}_{\text{KL}}[\pi_{\theta}(y|x) \mid \mid \pi^*(y\mid x)]$, with $\pi^*$ representing the optimal policy derived in Eq.~\ref{eq:optimum_model} for a reward function $r_{\phi}(y, x)$. With suitable algebraic manipulation, this yields the optimization problem:

\begin{equation}\label{eq:AC}
    \max_{\pi_{\theta}}\mathbb{E}_{\pi_{\theta}(y\mid x)}\bigg[\underbrace{r_{\phi}(x, y) -\beta\log\sum_{y}\pi_\text{ref}(y\mid x)\exp\left(\frac{1}{\beta}r_{\phi}(x, y)\right)}_{f(r_{\phi}, \pi_\text{ref}, \beta)} - \underbrace{\beta\log\frac{\pi_{\theta}(y\mid x)}{\pi_\text{ref}(y\mid x)}}_{\text{KL}}\bigg]
\end{equation}

This objective is identical to those targeted in previous studies 
\citep{ziegler2020finetuning, stiennon2022learning, bai2022training, ouyang2022training}, utilizing the DPO-equivalent reward corresponding to the reward function $r_{\phi}$. In this context, the normalization term within $f(r_{\phi}, \pi_\text{ref}, \beta)$ can be understood as representing the soft value function for the reference policy $\pi_\text{ref}$. Although this normalization component does not influence the solution that maximizes the objective, omitting it may result in a high-variance policy gradient, potentially destabilizing the training process. One way to address the normalization is by estimating it with a learned value function, though optimizing such a function can present its own challenges. A different approach taken in previous literature has been to normalize rewards with a baseline generated by human completions—effectively, a single-sample Monte Carlo approximation of the normalization constant. By contrast, the DPO reparameterization framework produces a reward function inherently free of the need for baselines.

\section{Experiments}
This section presents empirical results evaluating how effectively DPO can train policies based solely on preference data. We begin with a controlled text-generation experiment to compare DPO's efficiency in balancing reward maximization and KL-divergence minimization against standard preference optimization algorithms like PPO. Subsequently, we assess DPO's scalability and efficacy on larger-scale models and more complex RLHF applications, such as dialogue systems and summarization. Across these evaluations, we observe that DPO, with minimal hyperparameter adjustment, matches or outperforms robust benchmarks like RLHF with PPO and even the best-of-$N$ sampling approach using a learned reward model. Prior to reporting these findings, we outline the experimental protocol; further methodological specifics are provided in Appendix~\ref{app:exp_details}.

\textbf{Tasks.} In our study, we examine three distinct open-ended text generation tasks. Across all tasks, models are trained to optimize a policy based on a dataset of preferences, denoted as $\mathcal{D}=\bigl\{x^{(i)}, y_w^{(i)}, y_l^{(i)}\bigr\}_{i=1}^N$. The first task, \textbf{controlled sentiment generation}, uses $x$ as an initial segment from an IMDb movie review \cite{maas-EtAl:2011:ACL-HLT2011}, requiring the policy to generate a continuation $y$ expressing positive sentiment. For systematic evaluation in this setting, we automatically create pairs of candidate generations ranked by a pre-trained sentiment classifier, such that $p(\text{positive}\mid x,y_w)>p(\text{positive}\mid x,y_l)$. For the SFT baseline, GPT-2-large is fine-tuned to convergence using training reviews from the IMDb dataset (see App~\ref{app:sentiment_details} for further information). The second task, \textbf{summarization}, involves $x$ as a Reddit forum post, and the objective is for the policy to generate a summary $y$ covering its main points. Consistent with previous research, we utilize the Reddit TL;DR summarization dataset \citep{volske-etal-2017-tl}, along with the human preference annotations from \citeauthor{stiennon2022learning}. Our SFT baseline is obtained by fine-tuning on human-generated summaries of Reddit posts\footnote{\url{https://huggingface.co/CarperAI/openai_summarize_tldr_sft}} via the TRLX \citep{leandro_von_werra_2023_7790115} library for RLHF. Notably, the preference dataset was created by \citeauthor{stiennon2022learning} using samples from a distinct but similarly fine-tuned SFT model. The third task, \textbf{single-turn dialogue}, treats $x$ as a single user prompt, which can range widely from scientific inquiries to personal advice. Here, the policy must generate a response $y$ that is both helpful and engaging. We employ the Anthropic Helpful and Harmless dialogue dataset \citep{bai2022training}, which includes 170k conversations between a user and an automated assistant. At the end of each conversation, two responses generated by a large (but unspecified) language model are presented, together with a label indicating which the human raters preferred. In this context, as no pre-existing SFT model is accessible, we instead fine-tune a generic language model solely on preferred completions to construct the SFT model.

\textbf{Evaluation.} Our empirical studies incorporate two distinct evaluation methodologies. To assess the capacity of each algorithm to optimize the constrained reward maximization goal, we conduct analyses in the controlled sentiment generation framework by comparing their performance frontiers with respect to achieved reward and KL-divergence relative to the reference policy; this assessment is feasible since the ground-truth reward function—a sentiment classifier—is available. Conversely, in practical deployment scenarios where the true reward function is unknown, we utilize the metric of \textit{win rate} versus a baseline policy, employing GPT-4 as a surrogate for human judgment to evaluate summary quality and the helpfulness of responses in the summarization and single-turn dialogue tasks, respectively. For summarization, our baseline is constituted by reference summaries drawn from the test set, while for dialogue we compare against the response preferred in the test dataset. Though literature suggests that large language models outperform traditional automated metrics in evaluation roles \citep{Chen2023ExploringTU}, we perform a dedicated human study (described in Sec.~\ref{sec:human-judgments}) to support the suitability of using GPT-4 for evaluation. Our findings reveal that GPT-4’s judgments are highly correlated with those of humans, with the degree of agreement between human raters and GPT-4 matching or exceeding typical agreement among human annotators.

\textbf{Methods.} Alongside DPO, we assess a range of existing strategies for training language models to align with human preferences. As simple baselines, we investigate zero-shot prompting with \textbf{GPT-J} \citep{gpt-j} for summarization and two-shot prompting with \textbf{Pythia-2.8B} \citep{biderman2023pythia} in the dialogue domain. We further include the \textbf{SFT} model, as well as \textbf{Preferred-FT}, which involves supervised fine-tuning on the selected completion $y_w$ generated by either the SFT model (for sentiment control and summarization tasks) or a general-purpose language model (in the case of single-turn dialogue). Another supervised-style approach is \textbf{Unlikelihood}~\citep{welleck2019neural}, which explicitly maximizes the likelihood of $y_w$ while reducing the likelihood of $y_l$, employing an optional weighting factor $\alpha\in[0,1]$ for the unlikelihood term. We also test \textbf{PPO} \citep{schulman2017proximal}, leveraging a reward function inferred from preference annotations, as well as \textbf{PPO-GT}, an oracle agent utilizing the actual reward function accessible within the controlled sentiment context. In sentiment control experiments, we use two PPO-GT variants: an off-the-shelf implementation \cite{leandro_von_werra_2023_7790115} and a custom version incorporating reward normalization and advanced hyperparameter optimization for better results (these same adjustments are also applied to ‘standard’ PPO with learned reward functions). Lastly, we implement the \textbf{Best of $N$} baseline, where $N$ candidate completions are drawn from either the SFT model (or Preferred-FT for dialogue) and the candidate with the highest reward—assessed using a reward function trained on preference data—is selected. Despite its empirical strength, this approach separates reward modeling from PPO-based optimization, and is computationally inefficient at scale, as it necessitates generating $N$ completions per query during inference.

\subsection{How well can DPO optimize the RLHF objective?}

In standard RLHF approaches, the reward maximization is subject to a KL-divergence constraint, which serves to balance the incentive for higher reward against the need to keep the policy close to a reference policy. Consequently, any comparison between different algorithms should jointly consider both the achieved reward and the KL divergence from the reference; obtaining a marginal increase in reward at the expense of substantially larger KL divergence may be suboptimal. Figure~\ref{fig:frontier-tldr-main} depicts the tradeoff between reward and KL divergence—the reward-KL frontier—for several algorithms within the sentiment task context. For each algorithm, we conduct multiple runs, varying the policy conservativeness parameter in each run (target KL values of $\{3,6,9,12\}$ for PPO, $\beta$ values of $\{0.05,0.1,1,5\}$, $\alpha$ values of $\{0.05,0.1,0.5,1\}$ for unlikelihood, and using different random seeds for preferred-FT), totaling 22 experiments. Policies are evaluated on test prompts at intervals of 100 training steps, continuing until convergence. For each evaluation, we compute the average reward with respect to the true reward function and the mean sequence-level KL divergence\footnote{Calculated as the cumulative sum of KL-divergences at each timestep.} with the reference policy, $\text{KL}\left(\pi\mid \mid \pi_\text{ref}\right)$. The results demonstrate that DPO attains a considerably more favorable frontier than alternatives, achieving the highest reward while maintaining a low KL. Notably, this holds for several reasons. Although both DPO and PPO pursue identical optimization objectives, DPO achieves superior efficiency, delivering a more favorable reward-KL tradeoff that entirely dominates that of PPO. Furthermore, DPO attains a stronger frontier compared to PPO, \emph{even when PPO is provided access to ground-truth rewards} (PPO-GT).

\subsection{Can DPO scale to real preference datasets?}
\label{sec:dpo-real-datasets}
We proceed to assess the capability of DPO for fine-tuning on tasks involving summarization and single-turn dialogue. In the case of summarization, conventional automatic metrics like ROUGE often exhibit poor alignment with human judgments~\citep{stiennon2022learning}. Previous research has shown that utilizing PPO for fine-tuning language models in accordance with human preferences can yield more compelling summaries. Our comparison of approaches involves generating model outputs on the TL;DR summarization dataset’s test split, and then calculating the average frequency with which these outputs are preferred over reference summaries in the same set. For each method, outputs are sampled across a range of temperatures from 0.0 to 1.0, and the corresponding win rates are displayed in Figure~\ref{fig:frontier-tldr-main} (right). DPO, PPO, and Preferred-FT are all built on the same GPT-J SFT model\footnote{\url{https://huggingface.co/CarperAI/openai_summarize_tldr_sft}} for finetuning. Our experiments reveal that DPO achieves a win rate near 61\% at temperature 0.0, outperforming PPO, which attains about 57\% win rate at its own optimal setting (temperature 0.0). Additionally, the peak win rate for DPO surpasses that of the $N$-best baseline. It is important to highlight that the DPO $\beta$ hyperparameter received minimal tuning, suggesting the reported outcomes may not fully capture DPO's optimal performance. Furthermore, DPO maintains a high degree of resilience to changes in sampling temperature, unlike PPO, whose efficacy may drop to the level of the underlying GPT-J model at higher temperatures. Notably, Preferred-FT yields no substantial improvements over the initial SFT model. We extend our analysis with direct human comparison of DPO and PPO in Section~\ref{sec:human-judgments}, observing that DPO samples at temperature 0.25 are preferred in 58\% of cases compared to PPO samples at the same temperature.

For single-turn dialogue evaluation, we assess various approaches on the subset of the test split from the Anthropic HH dataset \citep{bai2022training} that contains a single human-assistant exchange. For GPT-4-based assessments, we treat the preferred test completions as references to determine win rates among the evaluated methods. Given the lack of an established SFT model for this domain, we initiate experiments with the pre-trained Pythia-2.8B, apply Preferred-FT on the selected completions to create a reference model whose outputs remain in-distribution, and subsequently employ DPO training. As baselines, we also consider the top-ranked completion among 128 Preferred-FT completions (since the Best of $N$ method shows diminishing returns at $N=128$ for this task; refer to Appendix Figure~\ref{fig:best-of-n}), as well as a 2-shot prompting variant using the Pythia-2.8B base. Our analysis indicates that DPO matches or surpasses baseline performance across the optimal temperature settings for each method. We further include in our comparison an RLHF model fine-tuned with PPO on the Anthropic HH dataset\footnote{\url{https://huggingface.co/reciprocate/ppo_hh_pythia-6B}} provided by a prominent source\footnote{\url{https://github.com/CarperAI/trlx/tree/main/examples/hh}}; however, we do not identify any prompting or sampling temperature that allows it to outperform the original Pythia-2.8B. Considering insights from our TL;DR experiments and the shared reward optimization objective, we treat Best of 128 as a loose estimate of PPO-level capability. In summary, DPO stands out as the only method that is both computationally efficient and capable of exceeding the preferred completion baseline on the Anthropic HH dataset, achieving performance on par with or better than the computationally intensive Best of 128 approach. Moreover, Figure~\ref{fig:dialogue-main} demonstrates that DPO rapidly approaches optimal results.

\subsection{Generalization to a new input distribution}

To better assess how PPO and DPO handle shifts in data distribution, we applied the policies learned from our Reddit TL;DR summarization task to a different data domain: news articles from the CNN/DailyMail test split \citep{nallapati-etal-2016-abstractive}. For evaluation, we utilized the same optimal sampling temperatures identified in the TL;DR experiment (0 and 0.25). Table~\ref{tab:ood} summarizes these results. Using the same GPT-4 (C) prompt as for Reddit TL;DR—substituting “forum post” with “news article”—we measured the GPT-4 win rate versus gold-standard summaries. On this out-of-domain benchmark, DPO again achieves substantially higher performance than PPO. These findings offer preliminary support that DPO’s generalization capabilities at least match those of PPO, despite DPO not incorporating the extra unlabeled Reddit TL;DR data employed by PPO.

\subsection{Validating GPT-4 judgments with human judgments}
\label{sec:human-judgments}
We carried out a human evaluation to test the robustness of GPT-4’s judgments by analyzing TL;DR summarization outputs under two different GPT-4 prompting strategies. The \textbf{GPT-4 (S)} prompt asks directly which summary best captures the core information of the post, while the \textbf{GPT-4 (C)} prompt additionally queries which summary is more concise, since we observed that GPT-4 using the (S) prompt favors lengthier, more repetitive responses compared to human preferences. Full prompt texts are provided in Appendix~\ref{app:prompts}. We conducted three distinct pairwise assessments to span a range of summary qualities, comparing the top (DPO, temp. 0.25), bottom (PPO, temp. 1.0), and an intermediate (SFT, temp. 0.25) approach against PPO with greedy decoding (its best-performing temperature). The analysis shows that, across both prompts, GPT-4’s choices correspond to human ratings with a frequency similar to inter-human agreement rates, implying GPT-4 is a reliable surrogate for human evaluation (note that only DPO and PPO-1 pairs received multiple independent human assessments, given our limited rater pool). Notably, the \textbf{GPT-4 (C)} prompt usually yields win rates more in line with human opinion; thus, this prompt is adopted for the primary results in Section~\ref{sec:dpo-real-datasets}. Further methodological details, including a description of the rating interface and information about participating human raters, can be found in Appendix~\ref{app:human-study}.

\section{Discussion}
Preference-based learning offers an effective and scalable approach for developing competent and aligned language models. In this work, we present DPO, a straightforward training approach that leverages human preferences to train language models without relying on reinforcement learning techniques. Rather than fitting preference learning into a traditional RL paradigm just to utilize standard RL algorithms, DPO establishes a correspondence between reward functions and language model policies. This allows the model to align with human preferences \textit{directly} through a simple cross-entropy loss, bypassing the need for reinforcement learning while preserving general applicability. With minimal hyperparameter tuning, DPO achieves performance on par with, or superior to, leading RLHF methods, including those based on PPO, thus significantly lowering the barrier for preference-based language model training.

\textbf{Limitations \& Future Work.} Our findings open up several avenues for future exploration. One central question is how well DPO-trained policies generalize outside the distribution, especially when contrasted with approaches trained on explicit reward signals. Preliminary evidence indicates that DPO policies generalize comparably to those trained with PPO, but a broader investigation is warranted. Another interesting direction is whether self-labeling strategies that utilize the DPO policy can take full advantage of prompts lacking human annotation. Furthermore, questions remain about how reward over-optimization might present itself within the direct preference optimization context, and whether the modest drop in performance depicted in Figure~\ref{fig:dialogue-main}-right exemplifies this phenomenon. While our evaluation currently extends to models with up to 6B parameters, further research into scaling DPO for much larger, state-of-the-art models promises substantial potential. With respect to evaluation, we observe that GPT-4’s calculated win rates can vary depending on the prompt used; thus, subsequent work could investigate optimal methodologies for obtaining robust assessments from automated evaluators. Lastly, DPO’s utility likely extends far beyond training language models with human preferences, encompassing a broad range of generative model training scenarios across different domains.